\chapter{知识图谱构建与管理}
\chapterart{ch07}{从原始数据到知识图谱的流水线}

\begin{introbox}
\textbf{【导读】}
前几章教你雕琢本体——一副精心打磨的概念骨架；但骨架自己不会回答业务问题，
真正承载价值的是骨架上的血肉：成千上万的实例与关系。
本章回答「血肉从哪里来、怎样长得健康」：
先盘点企业数据源，论证「先结构化主干」的开工顺序；
再走通两条进料通道——台账类数据的映射式构建、维修日志类文本的抽取式构建；
然后处理规模化必然撞上的三道坎：同一台设备的多个名字（实体对齐）、
错误合并的语义传染（\texttt{sameAs} 陷阱）、脏数据混入库内（SHACL 校验）；
最后讨论图谱住在哪里（三元组库与属性图的选型）、
怎样长期健康运转（具名图、监控与漂移检测）。
全章延续制造车间场景，代码片段全部节选自仓库
\texttt{ch07-knowledge-graph/examples/} 下四个可运行、可对照阅读的文件。
\end{introbox}

\section{本体与知识图谱：骨架与血肉}

本体评审通过的那天下午，会议室的气氛常常很微妙。
投影上是打磨数月的类层次——设备、工序、物料、故障，公理齐整，推理器全绿；
按第 3 章的方法论走完了全周期，按第 4 章的语法写得无可挑剔。
然后业务方代表客气地举手：「所以……车间里那几百台设备，什么时候能在系统里查到？」
会场安静了几秒钟。因为本体库里此刻只有几十个类和十几个手工敲进去的示例个体，
而业务要的是一张能查、能算、能追溯的\textbf{知识图谱}（knowledge graph）。
从「本体建好了」到「图谱可用了」，中间隔着一整条工程流水线——
这条流水线就是本章的主角。

\subsection{TBox 与 ABox 的规模断层}

第 2 章把知识库切成两层：TBox 是术语层，存放类、属性与公理；
ABox 是断言层，存放个体与事实。当时这是一个逻辑上的区分；
到了生产环境，它变成一道触目惊心的\textbf{规模断层}：

\begin{itemize}
\item \textbf{本体（TBox）}：数十到数百个概念，由本体工程师与领域专家逐条评审、
  人工精修；变更走发布流程，节奏以周、月计；
  每一条公理都有人能说出「为什么这么写」。
\item \textbf{知识图谱（TBox + ABox）}：百万到十亿条三元组，由自动化管道批量生产，
  节奏以分钟、小时计；没有任何人「看过每一条」，也永远不会有。
\end{itemize}

断层两侧是两个工程世界，治理手段截然不同。
骨架一侧靠\textbf{评审}：OntoClean 检查（第 3 章）、一致性检查（第 5 章）、同行评议，
质量靠「每条都被人看过」来保证。
血肉一侧靠\textbf{管道}：既然没人能看过每条三元组，
质量就只能由管道上的机器关卡保证——映射规则保证转换忠实，
实体对齐保证不重不漏，校验闸门把脏数据挡在库外，监控指标盯住整体健康度。
这也解释了一个常见误区：不少团队尝试在 Protégé 里手工录入实例，
几十台设备还录得动，几千条维修记录就彻底没指望了；
更糟的是手工录入既无溯源也无校验，错了都不知道从哪查起。
\textbf{实例数据永远走管道，手只用来写映射规则和裁决疑难案例}——这是本章的第一条纪律。

\subsection{知识图谱的工作定义}

「知识图谱」这个词在业界被用得很松，本书给它一个工程上可操作的工作定义：

\begin{noteblock}
\textbf{知识图谱 = 本体 + 实例 + 治理。}
本体提供词汇与公理（骨架），实例承载业务事实（血肉），
治理则是让两者持续健康咬合的全部机制——构建流水线、实体对齐、校验闸门、
溯源记录、监控与版本管理。三者缺一不可：
缺本体，图谱退化为字段语义不明的普通图数据；
缺实例，本体只是无血肉的标本；
缺治理，图谱会在几个月内淤积成无人敢信的「三元组沼泽」。
\end{noteblock}

前两个成分是第 2 章概念的直接延续——并且要注意，本体并未在图谱中「消失」：
\(\sqsubseteq\) 层次、不相交公理、函数性声明都随图谱一起入库并继续生效，
查询时沿子类闭包展开（7.8 节）、校验时充当判错依据（7.7 节）。
真正容易被遗漏的是第三个成分。许多项目把知识图谱当成「装满三元组的库」，
上线即巅峰，几个月后没人说得清某个值从哪来、对不对、新数据怎么进——
这不是技术失败，而是治理缺位。本章 7.5 节之后的全部篇幅，本质上都在讲治理。

\subsection{流水线总览与三道闸门}

把治理机制按数据流向串起来，就是贯穿本章的\textbf{构建流水线}：

\begin{center}
数据接入 \(\rightarrow\) 映射/抽取 \(\rightarrow\) 实体对齐 \(\rightarrow\)
本体校验闸门 \(\rightarrow\) 入库 \(\rightarrow\) 增量更新
\end{center}

六个环节各司其职：接入环节对接数据库与文档库（结构化源订阅变更、文档走批量解析）；
映射与抽取环节把原始记录变成候选三元组（7.3、7.4 节）；
对齐环节确保同一实体在图谱中只有一个节点（7.5、7.6 节）；
闸门拦下不合格批次（7.7 节）；入库环节落到选好的存储（7.8 节）；
增量环节让图谱跟上现实变化（7.9 节）。
其中\textbf{本体校验闸门}是质量生命线，值得在总览阶段就把它的三道检查立起来：

\begin{itemize}
\item \textbf{词汇检查}：候选三元组用到的类与属性必须已在本体中定义。
  抽取程序从日志里冒出一个 \texttt{mfg:SuperLathe}（「超导车床」），
  本体里查无此类——拒绝入库，转人工复核队列。
  这正是第 1 章预告过的「词汇硬约束」在流水线上的落地。
\item \textbf{公理检查}：候选事实不得与 TBox 公理冲突。
  同一台设备到货两个序列号，违反函数性公理；
  某设备同时被断言「运行中」与「维护中」，撞上不相交公理
  （运行中 \(\sqcap\) 维护中 \(\sqsubseteq\) \(\bot\)）——一律拒绝。
\item \textbf{形状检查}：SHACL 形状校验必填项、取值范围、格式与跨字段规则，
  7.7 节专门展开。
\end{itemize}

通过者写入，未通过者进\textbf{隔离区}并触发告警，由数据管理员裁决三选一：
数据错了（修源头）、抽取错了（修管道），还是本体漏了（走本体变更流程）。
这个裁决记录日后会成为知识漂移检测的信号源（7.9 节）。
闸门思想还会在第 9 章再次出场——那里把「先校验、后写入」做进系统实现的代码骨架，
本章先把每道关卡的原理与运营讲透。
下面按流水线顺序开讲，第一站：盘点你手里有什么数据。

\section{数据源盘点与优先级}

项目启动会上最常见的一幕：各系统负责人挨个报家底——
MES 数据库、台账 Excel、文档库里的工艺规程 PDF、维修班组的交接记录本，
白板上很快列出一长串数据源。「全部接进来」听起来雄心勃勃，
工程上却是灾难的开始：每类数据的进料方式、成本与回报差异巨大，
\textbf{接入顺序本身就是一个需要论证的设计决策}。
盘点的目的不是列清单，而是排优先级。

\subsection{三类数据源细看}

仓库文件 \texttt{kg-construction-pipeline.txt} 开篇就是盘点清单，按结构化程度分三档：

\begin{codebox}
\codeline{\textcolor{cmtgray}{\#\ ==============================}}
\codeline{\textcolor{cmtgray}{\#\ 1.\ 数据源盘点}}
\codeline{\textcolor{cmtgray}{\#\ ==============================}}
\codeline{\textcolor{cmtgray}{\#\ 结构化数据（最容易，优先做）：}}
\codeline{\textcolor{cmtgray}{\#\ \ \ 设备台账（Excel/MES数据库）——\ 设备、型号、功率、位置}}
\codeline{\textcolor{cmtgray}{\#\ \ \ BOM/工艺路线表\ \ \ \ \ \ \ \ \ \ \ \ \ ——\ 产品、工序、先后关系}}
\codeline{\textcolor{cmtgray}{\#\ 半结构化数据：}}
\codeline{\textcolor{cmtgray}{\#\ \ \ 工艺规程文件（Word/PDF表格）——\ 参数、约束条款}}
\codeline{\textcolor{cmtgray}{\#\ 非结构化数据（最难，价值高）：}}
\codeline{\textcolor{cmtgray}{\#\ \ \ 维修日志、班组交接记录\ \ \ \ \ \ ——\ 故障现象、原因、处置}}
\codeblank{}
\end{codebox}

\snipsource{ch07-knowledge-graph/examples/kg-construction-pipeline.txt}

\textbf{结构化数据}是图谱的「主粮」。
设备台账给出骨干实体：每台设备的编号、型号、功率、位置——
图谱里的设备节点几乎全部源于此；
BOM 与工艺路线表给出关系主干：产品由哪些工序加工、工序之间谁先谁后，
把孤立的设备节点织成一张网络。
这类数据字段语义清楚、有主键、有数据库接口，
映射规则写一次就能反复运行，是性价比最高的进料口。

\textbf{半结构化数据}介于两者之间。
工艺规程文件（Word/PDF）里藏着大量参数与约束条款——
主轴转速范围、进给量上限、特定材料的工艺禁忌。
表格解析器能拿到行列结构，但表头语义要人来对应、
合并单元格与跨页表格要逐模板适配；
一份模板调好之后批量处理同模板文件很顺，换一个模板又得再来一轮。

\textbf{非结构化数据}最难也最值钱。
维修日志与交接记录承载的是\textbf{经验知识}：什么现象、什么原因、怎么处置——
故障诊断与根因分析的原料全在这里，而这恰恰是台账给不了的。
代价是技术难度：自然语言、口语化表达、随手缩写、偶发错别字，
没有哪条正则能一劳永逸，7.4 节专门对付它。

\subsection{评估矩阵：把直觉变成表格}

盘点会的产出建议固化成一张评估矩阵，逐源打分。
打分不必精确，高/中/低三档足够支撑排序讨论：

\begin{center}
\begin{tabularx}{\linewidth}{@{}p{0.22\linewidth}XXXX@{}}
\toprule
\textbf{数据源} & \textbf{知识价值} & \textbf{数据质量} & \textbf{获取成本} & \textbf{更新频率} \\
\midrule
设备台账 & 高：骨干实体之源 & 高：有主键、字段规整 & 低：数据库直连或导出 & 低频，可订阅变更 \\
BOM/工艺路线 & 高：关系主干 & 较高：结构清晰 & 较低：多有系统接口 & 随产品换型批次更新 \\
工艺规程文件 & 中：参数与约束条款 & 中：模板不一、需解析 & 中：逐模板适配 & 低频，版本化发布 \\
维修日志、交接记录 & 高：故障经验知识 & 低：口语、缩写、错字 & 高：抽取加人工复核 & 持续新增 \\
\bottomrule
\end{tabularx}
\end{center}

怎么读这张矩阵？优先寻找「价值高、成本低」的西北角——台账与 BOM 恰好坐在那里。
这不是巧合：企业信息化多年沉淀下来的结构化系统，
本来就是按「最重要的业务对象先建系统」的逻辑生长的，
图谱建设顺着这个沉淀走，事半功倍。
维修日志价值同样高，但成本与质量两列把它推到第二梯队：
不是不做，而是\textbf{后做、小步做}。

另一个常被忽略的列是\textbf{更新频率}，它决定接入方式而非接入与否：
台账低频变更，适合 CDC（Change Data Capture，变更数据捕获）订阅增量；
维修日志持续新增，适合按日批量抽取；
工艺规程版本化发布，适合用发布事件触发解析。
接入方式在盘点阶段想清楚，7.9 节的增量运维才有地基。

\subsection{先结构化主干：一个值得背下来的开工顺序}

仓库文件末尾那条工程经验值得抄在项目白板上：
\textbf{先打通「台账到图谱」的结构化主干（两周可见效），再逐步纳入非结构化抽取}。
论证有三条，分别关于信任、依赖与风险。

\textbf{其一，见效快，建立信任。}
映射式构建是确定性工作：规则写对，数据就对。
两周之内业务方第一次在同一个界面里查到跨系统的设备全貌——
「全厂设备一张图」对管理层的说服力，胜过任何技术汇报；
项目从「研究课题」转成「生产工程」，往往就在这第一批可见产出上。

\textbf{其二，主干是抽取的锚点。}
维修日志抽出来的故障事件终究要挂到具体设备上
（\texttt{occursOn} 指向某台车床）；设备节点若不存在，
抽取结果就是无处着陆的碎片。
实体对齐（7.5 节）同样要靠台账提供的权威属性——型号、功率、位置——作比对基准。
先有主干，后来的知识才有挂靠之处。

\textbf{其三，风险后置。}
NLP 抽取的精度没有保底也没有速成通道，一开始就扎进去，
项目可能在长期没有可见产出的状态下耗尽预算与耐心。
把高确定性的映射工作放在前面，项目在任何时点都有东西可交付；
等主干稳定、团队对本体与管道都有了手感，再去啃难啃的文本。
「先结构化主干」的本质，是把不确定性最大的环节排到风险预算最充裕的阶段去。

盘点完成的标志是两件交付物：
一张数据源清单（负责人、接口方式、主键、更新机制、安全约束），
一个排好序的接入计划。下一节进入第一条进料通道。

\section{结构化数据映射：从表格到三元组}

第一条管道开工。手里是设备科维护多年的台账表，目标是本体约束下的三元组——
\textbf{映射式构建}（mapping-based construction）的全部秘密是一句话：
\textbf{规则一次写好，数据反复流过}。
W3C 为「关系数据库到 RDF」定过标准映射语言 R2RML（RDB to RDF Mapping Language）；
工程上也常见自写脚本，但无论用不用标准语法，R2RML 背后的映射观都值得先内化。

\subsection{R2RML 思想：四句口诀}

关系世界与图世界之间存在一组稳定的结构对应，记成四句口诀：

\begin{itemize}
\item \textbf{行 \(\rightarrow\) 个体}：表里一行就是一个实体，IRI 由稳定主键铸造；
\item \textbf{列 \(\rightarrow\) 数据属性}：普通列变成带类型标注的字面量断言；
\item \textbf{外键列 \(\rightarrow\) 对象属性}：指向另一张表主键的列，变成连接两个个体的边；
\item \textbf{代码表 \(\rightarrow\) 受控词表}：枚举值升格为本体中的命名个体。
\end{itemize}

对照仓库文件里的台账例子逐条验证：

\begin{codebox}
\codeline{\textcolor{cmtgray}{\#\ ==============================}}
\codeline{\textcolor{cmtgray}{\#\ 3.\ 结构化数据：映射式构建}}
\codeline{\textcolor{cmtgray}{\#\ ==============================}}
\codeline{\textcolor{cmtgray}{\#\ 台账表\ equipment\_ledger：}}
\codeline{\textcolor{cmtgray}{\#\ \ \ id\ \ \ \ |\ name\ \ \ \ \ \ |\ model\ \ \ |\ power\_kw\ |\ workcell}}
\codeline{\textcolor{cmtgray}{\#\ \ \ E0117\ |\ 3号数控车床\ |\ CK6140\ \ |\ 15.5\ \ \ \ \ |\ A区}}
\codeblank{}
\codeline{\textcolor{cmtgray}{\#\ R2RML/自写脚本映射为三元组：}}
\codeline{\textcolor{cmtgray}{\#\ \ \ mfg:Lathe\_003\ \ rdf:type\ \ \ \ \ \ \ \ \ mfg:CNCLathe\ ;}}
\codeline{\textcolor{cmtgray}{\#\ \ \ \ \ \ \ \ \ \ \ \ \ \ \ \ \ \ mfg:hasSerialNumber\ "EQ-2023-0117"\ ;}}
\codeline{\textcolor{cmtgray}{\#\ \ \ \ \ \ \ \ \ \ \ \ \ \ \ \ \ \ mfg:hasPower\ \ \ \ \ "15.5"\^{}\^{}xsd:float\ ;}}
\codeline{\textcolor{cmtgray}{\#\ \ \ \ \ \ \ \ \ \ \ \ \ \ \ \ \ \ mfg:locatedIn\ \ \ \ mfg:WorkCell\_A\ .}}
\codeblank{}
\codeline{\textcolor{cmtgray}{\#\ 映射要点：}}
\codeline{\textcolor{cmtgray}{\#\ \ \ 行\ \(\rightarrow\)\ 个体（IRI由稳定主键生成，禁止用易变的名称做IRI）}}
\codeline{\textcolor{cmtgray}{\#\ \ \ 列\ \(\rightarrow\)\ 数据属性；外键列\ \(\rightarrow\)\ 对象属性}}
\codeline{\textcolor{cmtgray}{\#\ \ \ 代码表（如状态枚举）\(\rightarrow\)\ 受控词表个体}}
\codeblank{}
\end{codebox}
\color{black}

\snipsource{ch07-knowledge-graph/examples/kg-construction-pipeline.txt}

一行台账（主键 E0117 的「3 号数控车床」，型号 CK6140）映射出四条三元组，每条都有讲究。
第一条 \texttt{rdf:type} 把个体挂进本体类 \texttt{mfg:CNCLathe}——
注意映射规则把型号家族对应到了具体的「数控车床」类，而不是泛泛的 \texttt{mfg:Equipment}：
归类越具体，第 5 章的推理服务能替你导出的结论越多。
第二条把序列号写成字符串字面量；台账主键 E0117 与全厂统一序列号 EQ-2023-0117
的换算就在映射规则里完成——编号归一做得越靠上游，7.5 节的对齐越省力。
第三条功率带上 \texttt{xsd:float} 类型标注；
第 4 章强调过「字面量要带类型」，在管道里这不是美德而是纪律：
下游的范围校验、数值比较、单位换算全部依赖类型。
第四条 \texttt{mfg:locatedIn} 由外键列 \texttt{workcell} 生成，
把设备连进车间拓扑，对象属性的值是另一个个体 \texttt{mfg:WorkCell\_A} 而非字符串「A区」——
\textbf{凡是会被继续追问的东西，都应当成为节点而不是字面量}。

\subsection{IRI 铸造：图谱里最早的不可逆决策}

每个个体的 IRI 在第一次入库时铸成，之后全图的边都压在它身上；
改 IRI 等于把这个实体的全部历史连根拔起再接回去，成本高到接近不可能。
所以 IRI 铸造原则必须在第一条管道上线之前定稿：

\begin{itemize}
\item \textbf{用稳定主键，禁用易变名称}。
  \texttt{mfg:Lathe\_003} 背后应当是台账主键或全厂序列号这类「死也不变」的标识。
  用名称铸 IRI 是经典事故源：「3 号数控车床」哪天改叫「3 号车铣复合中心」，
  IRI 跟着变，全图断链。仓库文件把这条写成了禁令——禁止用易变的名称做 IRI。
\item \textbf{命名空间规划先行}。本体词汇与实例数据共享
  \texttt{http://example.org/manufacturing\#} 这类命名空间时，
  前缀划分、大小写约定、编号格式都要立规矩；
  具名图另立 \texttt{http://example.org/graph/} 前缀（7.9 节）。
  命名空间一经发布就是对外承诺，宁可开工前多开一次评审会。
\item \textbf{IRI 不复用}。设备报废，IRI 保留并标注退役状态，
  绝不把旧 IRI 转给新设备——历史维修记录还挂在上面，
  复用等于把两台设备的一生缝在一起。
\end{itemize}

\subsection{代码表升格为受控词表}

台账里的状态字段是典型代码表：取值 Running/Idle/Maintenance/Fault。
映射时最常见的偷懒是把它们留作字符串字面量——
于是「running」「Running」「运行」三种写法并存，分组统计当场失真。
正确做法是按第四句口诀升格为\textbf{受控词表个体}：
\texttt{mfg:Running}、\texttt{mfg:Idle}、\texttt{mfg:Maintenance}、\texttt{mfg:Fault}，
状态断言变成对象属性指向这些个体。
好处链条很长：拼写歧义从机制上消失；
SHACL 用 \texttt{sh:in} 一行就能锁死取值集合（7.7 节会看到原文）；
词表个体还能继续挂属性——给 \texttt{mfg:Fault} 标上告警级别与处置预案，
所有故障设备的查询立即受益。
受控词表是「把约定从文档搬进数据」的最小实践，成本一次性，收益贯穿图谱终生。

\subsection{CDC 增量：让映射管道长跑}

全量重建只该发生一次——初始化；日常必须走增量，
否则每晚重灌百万级三元组的窗口迟早撑爆。
结构化源的增量靠 \textbf{CDC}：订阅台账数据库的变更流（插入、更新、删除），
把每条变更翻译成\textbf{三元组补丁}——一组 INSERT 与 DELETE 操作。
两个容易踩的坑：
其一，「更新」必须翻成「先删旧值、再插新值」，
三元组库没有「整行覆盖」的概念，漏删旧值就会出现一台设备两个功率并存；
其二，「删除」要区分语义——台账删行往往意味着设备报废而非录入笔误，
此时正确动作不是物理删除三元组，而是改挂退役标记，历史必须留下。
最后一条纪律：\textbf{补丁同样要过 7.7 节的闸门}。
增量是脏数据最爱走的侧门——全量时代验过的字段，
上游一次字段改版就会在增量补丁里再脏一次。
版本与回滚问题这里按下不表，7.9 节用具名图统一处理。

\section{非结构化抽取：从维修日志里挖知识}

主干跑通之后，轮到最难也最值钱的部分。先看一条真实形态的维修日志：

\begin{center}
「3月21日，3号车床主轴异响，检查发现轴承磨损，更换轴承后恢复运行。」
\end{center}

一句话里压着一台设备、一个部件、一种症状、一个原因、一次处置——
老师傅的诊断经验就藏在成千上万条这样的句子里，
台账永远不会告诉你「异响大概率是轴承磨损」。
把这些句子变成图谱里可查询、可统计、可推理的结构，
是\textbf{抽取式构建}（extraction-based construction）的任务。

\subsection{一条日志的完整解剖}

\begin{codebox}
\codeline{\textcolor{cmtgray}{\#\ ==============================}}
\codeline{\textcolor{cmtgray}{\#\ 4.\ 非结构化数据：抽取式构建}}
\codeline{\textcolor{cmtgray}{\#\ ==============================}}
\codeline{\textcolor{cmtgray}{\#\ 维修日志原文：}}
\codeline{\textcolor{cmtgray}{\#\ \ \ "3月21日，3号车床主轴异响，检查发现轴承磨损，更换轴承后恢复运行。"}}
\codeblank{}
\codeline{\textcolor{cmtgray}{\#\ 实体抽取（NER）：}}
\codeline{\textcolor{cmtgray}{\#\ \ \ [3号车床]\(\rightarrow\)Equipment\ \ [主轴]\(\rightarrow\)Component\ \ [异响]\(\rightarrow\)Symptom}}
\codeline{\textcolor{cmtgray}{\#\ \ \ [轴承磨损]\(\rightarrow\)FaultCause\ \ [更换轴承]\(\rightarrow\)RepairAction}}
\codeblank{}
\codeline{\textcolor{cmtgray}{\#\ 关系抽取（RE）：}}
\codeline{\textcolor{cmtgray}{\#\ \ \ (Lathe\_003,\ hasComponent,\ Spindle)}}
\codeline{\textcolor{cmtgray}{\#\ \ \ (FaultEvent\_0321,\ occursOn,\ Lathe\_003)}}
\codeline{\textcolor{cmtgray}{\#\ \ \ (FaultEvent\_0321,\ hasSymptom,\ AbnormalNoise)}}
\codeline{\textcolor{cmtgray}{\#\ \ \ (FaultEvent\_0321,\ hasCause,\ BearingWear)}}
\codeline{\textcolor{cmtgray}{\#\ \ \ (FaultEvent\_0321,\ resolvedBy,\ BearingReplacement)}}
\codeblank{}
\codeline{\textcolor{cmtgray}{\#\ 抽取技术选型：}}
\codeline{\textcolor{cmtgray}{\#\ \ \ 规则/模板\ \ ——\ 高精度低召回，适合格式固定的日志}}
\codeline{\textcolor{cmtgray}{\#\ \ \ NER模型\ \ \ \ ——\ 需标注数据，适合大批量同类文本}}
\codeline{\textcolor{cmtgray}{\#\ \ \ LLM抽取\ \ \ \ ——\ 零样本能力强，但必须经本体校验后入库（第8章详述）}}
\codeblank{}
\end{codebox}
\color{black}

\snipsource{ch07-knowledge-graph/examples/kg-construction-pipeline.txt}

抽取分两步。\textbf{实体抽取}（NER，Named Entity Recognition）把文本片段标到本体类上：
「3号车床」是 Equipment、「主轴」是 Component、「异响」是 Symptom、
「轴承磨损」是 FaultCause、「更换轴承」是 RepairAction。
注意标签集直接取自本体类，而不是 NLP 工具默认的「人名、地名、机构」三件套——
本体在这里第一次\textbf{反向}服务于 NLP 任务：它就是抽取的目标模式（schema）。

\textbf{关系抽取}（RE，Relation Extraction）把实体连成三元组。
最值得驻足的是其中的建模决策：故障这件「事」被铸成一个中介节点
\texttt{FaultEvent\_0321}，症状、原因、处置全挂在事件上，而不是直接挂在设备上。
为什么？同一台设备一生会发生很多次故障，
直接挂设备会把不同次故障的属性搅成一锅粥；
事件节点让每次故障保持独立、可计数、可排序、可关联时间与人员——
这正是第 4 章「多参与者关系用中介节点」模式的实战重演。
另外注意 \texttt{(Lathe\_003, hasComponent, Spindle)} 这类「顺手抽出」的结构知识：
日志顺带告诉了我们车床有主轴，这条边对后续的故障传播分析很有价值——
抽取管道的产出常常超出最初的需求清单，前提是本体里给这些知识留了位置。

\subsection{三代抽取技术的成本账}

仓库文件末尾的技术选型可以展开成一张成本对照表：

\begin{center}
\begin{tabularx}{\linewidth}{@{}p{0.13\linewidth}p{0.24\linewidth}XX@{}}
\toprule
\textbf{技术代际} & \textbf{前期投入} & \textbf{精度与召回特征} & \textbf{适用场景} \\
\midrule
规则/模板 & 写规则的人力，无需标注 & 高精度、低召回，漏掉规则未覆盖的表述 & 格式固定的日志，编号、日期等规整字段 \\
NER 模型 & 成规模的标注语料与训练调参 & 精度召回均衡，对标注分布内的文本稳定 & 大批量同类文本的长期管道 \\
LLM 抽取 & 几乎零样本起步，主要是提示词工程 & 召回强、覆盖灵活，但会产出本体之外的「幻觉」结果 & 文体多变、标注稀缺的冷启动，必须配本体校验闸门 \\
\bottomrule
\end{tabularx}
\end{center}

三代不是替代关系，而是同一工具箱里的三件工具：
规整字段（日期、编号、规格）用正则又快又稳，没必要劳驾模型；
量大且文体稳定的核心日志值得养一个 NER 模型，长期边际成本最低；
长尾文体与冷启动阶段交给 LLM，零样本能力无可替代。
真实管道几乎总是混合编队——规则打底、模型扛主力、LLM 补盲区，
按「每千条文本的综合处理成本」而不是「技术先进性」来分配地盘。

\subsection{LLM 抽取与本体校验的闭环}

LLM 把抽取的冷启动成本压到了前所未有的低——
不标注、不训练，把日志和「请按这些类与属性抽取」的提示词一起发过去，
就能拿到结构化结果。
但它有一个符号系统绝不会犯的毛病：\textbf{无中生有}。
日志里写「主轴异响」，它可能体贴地发明一个本体里不存在的设备类型
（7.1 节那个 \texttt{mfg:SuperLathe} 就是这么来的），
或者把没写明的功率「合理推断」出来。
工程答案不是把模型调到不出错（做不到），而是\textbf{闭环}：

\begin{center}
LLM 抽取 \(\rightarrow\) 词汇检查 \(\rightarrow\) 公理与形状检查 \(\rightarrow\)
通过入库；未通过 \(\rightarrow\) 带违规报告回到复核队列 \(\rightarrow\) 人工改判回流为示例
\end{center}

本体在闭环里扮演双重角色：抽取前它是\textbf{提示词里的词汇表}，约束输出空间；
抽取后它是\textbf{校验的判错依据}，兜住越界输出。
「生成靠神经网络、把关靠符号系统」——这套组合拳是第 8 章的主题，
那里会给出提示词模板、Text2SPARQL 与幻觉控制的完整方案；
本章的任务是先把闸门修好：闸门不牢，LLM 进来的就不是知识而是污染。

\subsection{抽取质量怎么评}

管道上线前必须回答「抽得准不准」，方法是\textbf{抽样建金标准}：
从日志里随机抽一批，人工标出全部实体与关系作为标准答案，
然后分槽位计算\textbf{精确率}（抽出来的有多少是对的）与\textbf{召回率}（该抽的抽出来多少）。
两个运营要点。
其一，\textbf{按关系类型分别报数}，整体平均会掩盖短板——
\texttt{hasSymptom} 往往很准（症状词汇集中），
\texttt{hasCause} 通常差一截（因果表述的语言变化远比症状丰富），
两者混在一个平均数里，短板就被藏起来了。
其二，\textbf{上线后持续抽检}而非一锤子验收：
日志文体会漂移——新班组、新设备、新缩写——精度会无声衰减，
7.9 节的漂移检测会回收这个话题。
最后是兜底原则：评估做得再好也有漏网之鱼，
所以文本抽取的三元组一律进单独的低可信具名图（7.9 节），与权威数据物理隔离——
\textbf{质量不确定的知识可以入库，但必须可识别、可隔离、可回滚}。

\section{实体对齐：同一台设备的多个名字}

第一次用图谱做全厂设备统计，业务方就皱起了眉：
图谱里的「设备总数」比台账多出一截。
排查发现不是多了设备，而是\textbf{同一台设备在图谱里裂成了多个节点}：
台账叫「3号数控车床」，MES 叫「CNC-Lathe-003」，
维修日志里写「3\#车床」「三号车床」甚至「L003」——
各来源各自映射、各自铸 IRI，互不相认。
设备数量虚高还只是难看，致命的是下游分析全面失真：
故障率被摊薄到几个分身头上，维修历史四分五裂，
「这台设备去年修过几次」这种最基本的问题都答不对。
\textbf{实体对齐}（entity resolution）因此是知识融合的第一关——
对齐不做好，前两节辛苦建的管道只是在更快地制造混乱。

\subsection{问题的两个方向}

\begin{codebox}
\codeline{\textcolor{cmtgray}{\#\ ==============================}}
\codeline{\textcolor{cmtgray}{\#\ 1.\ 问题形态}}
\codeline{\textcolor{cmtgray}{\#\ ==============================}}
\codeline{\textcolor{cmtgray}{\#\ 同一实体多种写法（需合并）：}}
\codeline{\textcolor{cmtgray}{\#\ \ \ 台账系统：\ \ "3号数控车床"\ \ (id=E0117)}}
\codeline{\textcolor{cmtgray}{\#\ \ \ MES系统：\ \ \ "CNC-Lathe-003"}}
\codeline{\textcolor{cmtgray}{\#\ \ \ 维修日志：\ \ "3\#车床"\ /\ "三号车床"\ /\ "L003"}}
\codeblank{}
\codeline{\textcolor{cmtgray}{\#\ 不同实体同名（需区分）：}}
\codeline{\textcolor{cmtgray}{\#\ \ \ "加工中心"\ ——\ A车间的\ VMC850\ 与\ B车间的\ VMC1060\ 都被这么叫}}
\codeblank{}
\end{codebox}

\snipsource{ch07-knowledge-graph/examples/entity-resolution.txt}

对齐问题有两个相反的方向，工程上都要防。
\textbf{同一实体多种写法}是主流形态，根源是各系统独立建设、各起各名，
解决方向是\textbf{合并}。
\textbf{不同实体同名}更阴险：A 车间的 VMC850 与 B 车间的 VMC1060
在口语里都叫「加工中心」，若按名称草率合并，
两台设备的维修史搅成一团，比不合并还糟——
解决方向是\textbf{区分}，靠名称之外的证据（车间、型号、功率）把它们撑开。

所有对齐技术都在两种错误之间走钢丝：
\textbf{漏合}（该合没合）让图谱碎片化，\textbf{错合}（不该合却合了）让图谱被污染。
关键认知是\textbf{两种代价不对称}：碎片还能在后续批次里继续缝合，
而被污染的历史很难洗净——错合的恶果还会沿着查询与推理继续扩散（7.6 节细讲机制）。
这个不对称将决定本节的阈值设计与下一节的语义选择，请先记住它。

\subsection{第一层：规则归一化}

工程上的对齐分两层，先规则后模型。规则层朴素得不像「技术」，
却承包了对齐工作量的大头：

\begin{codebox}
\codeline{\textcolor{cmtgray}{\#\ ==============================}}
\codeline{\textcolor{cmtgray}{\#\ 2.\ 第一层：规则归一化}}
\codeline{\textcolor{cmtgray}{\#\ ==============================}}
\codeline{\textcolor{cmtgray}{\#\ 编号正则归一：}}
\codeline{\textcolor{cmtgray}{\#\ \ \ "3号|3\#|三号|No.3"\ \ +\ "车床"\ \ \(\rightarrow\)\ \ 候选\ Lathe\_003}}
\codeline{\textcolor{cmtgray}{\#\ \ \ 设备编号\ "E0117"\ /\ "EQ-2023-0117"\ \(\rightarrow\)\ 去前缀补零归一}}
\codeblank{}
\codeline{\textcolor{cmtgray}{\#\ 别名表（最朴素也最有效）：}}
\codeline{\textcolor{cmtgray}{\#\ \ \ Lathe\_003:\ \ ["3号数控车床",\ "CNC-Lathe-003",\ "3\#车床",\ "L003"]}}
\codeline{\textcolor{cmtgray}{\#\ \ \ 维护成本低，命中即对齐；新别名出现时人工补一行}}
\codeblank{}
\codeline{\textcolor{cmtgray}{\#\ 规格归一化：}}
\codeline{\textcolor{cmtgray}{\#\ \ \ 功率\ "15.5kW"\ /\ "15500W"\ /\ "15.5千瓦"\ \(\rightarrow\)\ 15.5\ (kW,\ xsd:float)}}
\codeline{\textcolor{cmtgray}{\#\ \ \ 日期\ "2026/3/21"\ /\ "26年3月21日"\ \ \ \ \ \(\rightarrow\)\ "2026-03-21"\^{}\^{}xsd:date}}
\codeblank{}
\end{codebox}
\color{black}

\snipsource{ch07-knowledge-graph/examples/entity-resolution.txt}

\textbf{编号正则}处理写法变体：「3号、3\#、三号、No.3」加上设备类型词，
统一归到候选 \texttt{Lathe\_003}；
台账主键 E0117 与序列号 EQ-2023-0117 之间「去前缀、补零」即可互换。
编号体系的换算规则，设备科的老师傅一张嘴就能说全——
\textbf{先访谈再写规则}，比从数据里反推快得多。

\textbf{别名表}是性价比之王：一张「主 IRI 对别名列表」的映射表，
命中即对齐，零歧义、零延迟；新别名出现时人工补一行。
它看起来笨拙，常被技术团队跳过直接上模型——这是错误的。
算一笔账：别名表的维护成本随时间\textbf{递减}（高频别名很快收敛齐全），
模型的维护成本随时间\textbf{累积}（数据漂移就要重新训练调参）。
车间设备是一个\textbf{封闭且缓慢变化的实体集}——
几百台设备、每台几个别名，一张表就是全集。
这种场景下别名表不是权宜之计，而是正解；
模型适合的是开放实体集（比如不断冒出来的供应商与物料），别拿错工具。

\textbf{规格归一化}对付属性值的写法分裂：
功率「15.5kW」「15500W」「15.5千瓦」统一为 15.5 加 \texttt{xsd:float}（千瓦量纲）；
日期「2026/3/21」「26年3月21日」统一为 \texttt{xsd:date} 的「2026-03-21」。
注意这一步不只是为对齐服务：
下一小节的属性相似度、7.7 节的范围校验，
都依赖「同一量纲、同一类型」这个前提。归一化做得越靠上游，下游越干净。

\subsection{第二层：相似度匹配与双门限}

规则没接住的长尾进入第二层：给候选实体对打分，按分数分流。

\begin{codebox}
\codeline{\textcolor{cmtgray}{\#\ ==============================}}
\codeline{\textcolor{cmtgray}{\#\ 3.\ 第二层：相似度匹配}}
\codeline{\textcolor{cmtgray}{\#\ ==============================}}
\codeline{\textcolor{cmtgray}{\#\ 规则未命中时，计算候选对的加权相似度：}}
\codeline{\textcolor{cmtgray}{\#\ \ \ 名称相似度（编辑距离/拼音）\ \ \ \ \ \ \ \ \ \ 权重\ 0.3}}
\codeline{\textcolor{cmtgray}{\#\ \ \ 属性相似度（型号、功率、厂商一致性）\ \ 权重\ 0.4}}
\codeline{\textcolor{cmtgray}{\#\ \ \ 关系相似度（所在工作单元、操作工重合）权重\ 0.3}}
\codeblank{}
\codeline{\textcolor{cmtgray}{\#\ 判定阈值：}}
\codeline{\textcolor{cmtgray}{\#\ \ \ score\ \(\ge\)\ 0.90\ \ \(\rightarrow\)\ 自动合并}}
\codeline{\textcolor{cmtgray}{\#\ \ \ 0.70\ \textasciitilde{}\ 0.90\ \ \ \(\rightarrow\)\ 进人工复核队列}}
\codeline{\textcolor{cmtgray}{\#\ \ \ score\ <\ 0.70\ \ \(\rightarrow\)\ 视为不同实体}}
\codeblank{}
\codeline{\textcolor{cmtgray}{\#\ 示例：}}
\codeline{\textcolor{cmtgray}{\#\ \ \ ("3号车床",\ power=15.5,\ cell=A)\ vs\ (CNC-Lathe-003,\ power=15.5,\ cell=A)}}
\codeline{\textcolor{cmtgray}{\#\ \ \ 名称0.6\(\times\)0.3\ +\ 属性1.0\(\times\)0.4\ +\ 关系1.0\(\times\)0.3\ =\ 0.88\ \(\rightarrow\)\ 人工复核\ \(\rightarrow\)\ 确认合并}}
\codeblank{}
\end{codebox}

\snipsource{ch07-knowledge-graph/examples/entity-resolution.txt}

三路证据加权。\textbf{名称相似度}（编辑距离、拼音）权重 0.3，
\textbf{属性相似度}（型号、功率、厂商一致性）权重 0.4，
\textbf{关系相似度}（所在工作单元、操作工重合度）权重 0.3。
权重分配本身就是领域判断：名称权重最低，
因为名称恰恰是各系统分裂最严重的字段；
属性与关系合占七成——「叫什么」不可靠，
「长什么样、在哪儿、跟谁打交道」才可靠。
关系相似度是图谱独有的武器：普通主数据匹配只能比字段，
图谱还能比\textbf{邻居}，两个节点若连着同一个工作单元和同一批操作工，
是同一台设备的证据就很硬。

文件里的例子值得亲手算一遍：「3号车床」与「CNC-Lathe-003」字面上风马牛不相及，
名称相似度只有 0.6（拼音与编号成分救回一些分），
但功率同为 15.5、同在 A 区，属性与关系两路都是满分 1.0，
加权总分 0.6\(\times\)0.3 + 1.0\(\times\)0.4 + 1.0\(\times\)0.3 = 0.88。

接下来是分流。判定不用单一阈值，而是\textbf{双门限三段制}：
得分不低于 0.90 自动合并；0.70 至 0.90 之间进人工复核队列；
低于 0.70 判为不同实体。0.88 落在复核区，
复核员对照台账与 MES 的原始记录看一眼，确认合并。

为什么是双门限？单一阈值逼着你在错合与漏合之间二选一：
线划高了漏合暴增，线划低了错合渗入，而我们已经知道错合的代价高得多。
双门限承认「机器拿不准的就是拿不准」，把不确定的中间带显式交给人：
机器吞掉高分与低分两端的处理量大头，人只看中间窄带——
这是人机分工里最划算的切法。

复核队列是\textbf{运营}出来的，不是配置出来的：
队列要有专人值守与处理时限——积压意味着新数据延迟入库；
每个判例要留下「为什么合、为什么不合」的一句话理由，这是日后调参的原料；
仓库文件最后一条说得直白：阈值不是一次定终身，
\textbf{每月用复核队列攒下的人工标注做回归，重新校准权重与门限}。
人工复核不是成本，是免费送上门的标注数据，不回流就浪费了。

\subsection{上线前的对齐审计}

对齐质量直接决定图谱可信度，上线前要过\textbf{审计}：
随机抽取 200 个合并决策交人工复核，精确率应不低于 98\%。
两个数字都来自仓库文件，红线划在这里的理由值得说透：
合并是一种会\textbf{传播}的操作——错合一对实体，
它们名下的全部历史互相污染，并沿着后续的查询、统计、推理继续扩散。
98\% 意味着每一百个合并决策至多容两个错，
这是「污染型错误」场景下工程上可接受的水位；
若审计不过线，回到上一小节调权重、提门限、扩别名表，过线才放行。
还要注意审计对象是\textbf{合并决策}而非全体实体：
漏合不在这条红线里——它的代价是碎片化，可以靠后续批次继续缝合，紧迫性低一档。
审计与 7.4 节的抽取评估、7.7 节的形状校验合在一起，
构成图谱质量的三道独立防线：抽得准、认得对、格式合规。

\section{sameAs 的语义与联邦陷阱}

对齐判定「这两个节点是同一台设备」之后，还剩一个看似平淡的问题：
\textbf{在图谱里怎么表达这个结论}？
这一步选错语义，前面所有的小心翼翼都会泡汤——
因为表达方式决定了错误（总会有错误漏网）的爆炸半径。

\subsection{两种表达：物理合并与 sameAs}

\begin{codebox}
\codeline{\textcolor{cmtgray}{\#\ ==============================}}
\codeline{\textcolor{cmtgray}{\#\ 4.\ 合并的表达：owl:sameAs\ 及其陷阱}}
\codeline{\textcolor{cmtgray}{\#\ ==============================}}
\codeline{\textcolor{cmtgray}{\#\ 确认对齐后的两种做法：}}
\codeblank{}
\codeline{\textcolor{cmtgray}{\#\ 做法A：物理合并（推荐用于同一企业内部）}}
\codeline{\textcolor{cmtgray}{\#\ \ \ 选定主IRI（mfg:Lathe\_003），其余标识写成别名属性：}}
\codeline{\textcolor{cmtgray}{\#\ \ \ mfg:Lathe\_003\ mfg:hasAlias\ "CNC-Lathe-003",\ "3\#车床"\ .}}
\codeblank{}
\codeline{\textcolor{cmtgray}{\#\ 做法B：owl:sameAs\ 关联（用于跨系统联邦场景）}}
\codeline{\textcolor{cmtgray}{\#\ \ \ mes:CNC-Lathe-003\ owl:sameAs\ mfg:Lathe\_003\ .}}
\codeline{\textcolor{cmtgray}{\#\ \ \ 陷阱：sameAs是强语义——两个IRI的所有断言互相传播；}}
\codeline{\textcolor{cmtgray}{\#\ \ \ 若对齐错误，错误属性会"传染"到正确实体。}}
\codeline{\textcolor{cmtgray}{\#\ \ \ 联邦场景中宁可用弱关联\ skos:closeMatch，确认后再升级为sameAs}}
\codeblank{}
\end{codebox}

\snipsource{ch07-knowledge-graph/examples/entity-resolution.txt}

做法 A 是\textbf{物理合并}：选定主 IRI（\texttt{mfg:Lathe\_003}），
把其余标识降格为别名属性 \texttt{mfg:hasAlias}，
所有来源的断言改写到主 IRI 名下。
图谱里始终只有一个节点，查询、统计、推理都干净利落；
代价是合并动作要改写数据，撤销时要靠溯源记录倒放。
\textbf{同一企业内部建议默认走这条路}——既然在同一个治理域里，
就没有理由让一台设备保留多个分身。

做法 B 保留两个 IRI，用 \texttt{owl:sameAs} 桥接，适用于\textbf{跨系统联邦}：
MES 的数据归 MES 团队持续维护、按自己的节奏发布，
你无权也不应改写对方的库，于是各自保留节点、声明等同。
听起来温和，实际上 \texttt{owl:sameAs} 是 OWL 里语义最烈的词汇之一，
值得专门用一小节讲清它的脾气。

\subsection{传播机制与错误传染}

\texttt{owl:sameAs} 的语义是\textbf{个体等同}：两个 IRI 指称同一个对象。
推理器随即把它们当作同一个节点对待——
A 名下的每条断言都同样成立于 B，反之亦然，\textbf{全部断言双向传播}。
这正是它好用的地方：声明一条 sameAs，两个系统的知识自动汇流，
查任何一个 IRI 都能看到完整画像。
也正是它可怕的地方：传播不辨良莠，错误同样全速扩散。

拿 7.5 节那对同名设备做沙盘推演。
某次草率对齐把 A 车间的 VMC850 与 B 车间的 VMC1060 连上了 sameAs——
两台「加工中心」，名称相似度确实高。推理器随即开工：
VMC850 的维修记录传播到 VMC1060 名下，B 车间看到一堆自己从没做过的维修史；
两台设备的功率互相传播，每台名下出现两个不同的功率值——
若功率被声明为函数性属性，第 5 章讲过的一致性检查直接亮红灯，
而解释服务给出的最小公理集会指向那条无辜的功率断言，元凶 sameAs 反而藏在解释链的深处；
位置断言传播之后，一台设备同时「位于」两个车间，排产与库存系统跟着混乱。
最阴险的是\textbf{清理难}：若推理结论已经物化进库（第 5 章的物化策略），
撤销那条 sameAs 并不会自动收回衍生三元组，
必须靠溯源记录把每一条传播产物找出来删掉。
没有溯源的图谱遇到这种事故，等于做一次没有备份的灾难恢复。

\subsection{closeMatch 升级阶梯}

联邦场景的工程答案，是把「等同」从一锤子买卖改成\textbf{升级阶梯}：

\begin{center}
对齐候选 \(\rightarrow\) \texttt{skos:closeMatch}（弱关联，不传播）
\(\rightarrow\) 人工确认或审计通过 \(\rightarrow\) \texttt{owl:sameAs} 或物理合并
\end{center}

\texttt{skos:closeMatch} 来自 SKOS 词汇表，语义只是「两者非常接近」：
人看了知道有关联，应用可以按需利用，但\textbf{不触发任何断言传播}——
推理器不会替你做任何危险的事。
对齐算法产出的候选一律先停在这一档；
经人工确认或抽样审计达标后，再升级为 \texttt{owl:sameAs}（或回到做法 A 物理合并）；
拿不准的就让它一直停着——一条弱关联放着不动没有任何代价，
一条错误的 sameAs 每分钟都在传染。
升级动作本身也要走流程：谁确认的、依据什么证据、何时升级，记入溯源。
这套阶梯与 7.5 节的双门限一脉相承：
\textbf{把确定性分级，强语义只授予经过验证的结论}。

\subsection{对齐之后：属性冲突的消解}

合并（或桥接）完成，立刻浮出下一个问题：
两个来源对同一属性各执一词怎么办？

\begin{codebox}
\codeline{\textcolor{cmtgray}{\#\ ==============================}}
\codeline{\textcolor{cmtgray}{\#\ 5.\ 冲突消解}}
\codeline{\textcolor{cmtgray}{\#\ ==============================}}
\codeline{\textcolor{cmtgray}{\#\ 对齐后属性值冲突：台账说功率15.5，MES说11.0}}
\codeline{\textcolor{cmtgray}{\#\ 消解策略（按优先级）：}}
\codeline{\textcolor{cmtgray}{\#\ \ \ (1)\ 来源可信度：台账(设备科维护)\ >\ MES配置\ >\ 日志文本抽取}}
\codeline{\textcolor{cmtgray}{\#\ \ \ (2)\ 时效性：带时间戳的新值覆盖旧值（保留历史到具名图）}}
\codeline{\textcolor{cmtgray}{\#\ \ \ (3)\ 提交裁决：高价值字段（安全相关）冲突必须人工确认}}
\codeblank{}
\codeline{\textcolor{cmtgray}{\#\ 溯源记录（与第5章PROV思想一致）：}}
\codeline{\textcolor{cmtgray}{\#\ \ \ 每条三元组附带来源与时间元数据，}}
\codeline{\textcolor{cmtgray}{\#\ \ \ 冲突裁决时可追溯"这个值是谁、什么时候、从哪个系统来的"}}
\codeblank{}
\end{codebox}

\snipsource{ch07-knowledge-graph/examples/entity-resolution.txt}

台账说功率 15.5，MES 说 11.0，图谱里留哪个？消解策略按优先级排队。
\textbf{第一优先：来源可信度}。台账由设备科专人维护，是设备属性的权威源；
MES 配置次之；日志文本抽取的值垫底。
落地方式是给每个来源（具名图）赋可信度等级，冲突时高级别胜出——
这要求 7.9 节的按来源分图先行，否则连「这个值从哪来」都答不出。
\textbf{第二优先：时效性}。带时间戳的新值覆盖旧值，
旧值不删除而是退入历史具名图——设备确实会改造升级，
「冲突」有时只是「变化」，时间戳能把两者区分开。
\textbf{第三优先：人工裁决}。安全相关的高价值字段（额定功率、承重上限这类），
任何冲突都必须人来确认，机器只负责把冲突摆上桌面并附全证据。

三条策略能够运转的共同前提是\textbf{溯源}（provenance）：
每条三元组都带着「谁、什么时候、从哪个系统来」的元数据——
这正是第 5 章 PROV 思想与具名图机制在管道里的用武之地。
没有溯源，冲突消解退化成「最后写入者赢」，
而那是数据治理里最坏的策略：结果取决于管道调度的偶然顺序。

\section{SHACL 数据质量工程}

7.1 节给校验闸门立了三道检查：词汇检查与公理检查依托本体自身即可执行，
真正吃掉日常工作量的是第三道——形状检查。
先看几条真实管道里被拦下来的「货」：
一台设备没有序列号（台账早年的记录漏填）；
一台设备功率是 0（Excel 单元格的默认值混了进来）；
一条状态断言写着「检修」（上游新同事没用受控词表）；
一条维护记录的维护日期晚于记录创建日期（导出时日期列错位）。
难堪之处在于，这些数据在 RDF 语法上全部无懈可击——
三元组模型百无禁忌，任何主谓宾组合都是合法 RDF；
可它们一旦入库，下游的查询、统计与推理都会跟着失真。
关系数据库时代，这道防线由表结构与应用层把守
（非空列、检查约束、表单校验）；图谱时代，
它要用 \textbf{SHACL}（Shapes Constraint Language，形状约束语言）重建。

\subsection{OWL 推理为什么管不了数据质量}

第一反应常常是：本体里不是写了公理吗？
给序列号属性声明「最少一个」的基数公理，缺序列号不就会被推理器揪出来？
偏偏不会。第 2 章埋下的开放世界假设在这里结出工程后果：
推理器面对「图谱里查不到这台设备的序列号」，解读是
「序列号存在，只是尚未被记录」——基数公理甚至会帮它推出
「存在一个匿名的序列号」这种结论，全程心安理得，毫无报错。
仓库文件 \texttt{kg-quality-shacl.ttl} 开头的注释一语中的：
\textbf{缺序列号在 OWL 下只是「未知」，在 SHACL 下就是一条违规报告}。
把两套机制的脾气并排放清楚：

\begin{center}
\begin{tabularx}{\linewidth}{@{}p{0.16\linewidth}XX@{}}
\toprule
\textbf{维度} & \textbf{OWL 公理（推理）} & \textbf{SHACL 形状（校验）} \\
\midrule
世界假设 & 开放世界：没说不等于没有 & 封闭检查：图里没有就是违规 \\
缺失必填项 & 沉默，甚至推出匿名个体补位 & 一条违规报告，定位到具体节点 \\
职责 & 导出新知识、检出逻辑矛盾（第 5 章） & 把不合格数据挡在库外 \\
归属与节奏 & TBox 的一部分，随领域语义演化 & 独立形状文件，随质量需求演化 \\
出错形态 & 整库不一致，定位靠解释服务 & 逐节点违规清单，自带提示消息 \\
\bottomrule
\end{tabularx}
\end{center}

结论是分工而非取舍：公理负责语义，说清这个世界「是什么」；
形状负责质量，检查这批数据「合不合格」。
把 OWL 公理当校验器，是新手最常见的误用——公理越写越多，
推理越来越慢、意外结论越来越多，脏数据却一条没拦住；
反过来想用 SHACL 顶替本体也不行，形状不会替你做任何第 5 章的推理服务。
记一句口诀：\textbf{推理向外扩（导出没说的），校验向内收（拒绝不该有的）}。

\subsection{一个形状的解剖}

SHACL 自己也用 RDF 书写，第 4 章的 Turtle 语法原样复用，不需要任何新工具链。
核心概念只有两个：\textbf{节点形状}（node shape）圈定「哪些节点受检」，
其下的若干\textbf{属性形状}（property shape）规定「每个属性要满足什么」。
看设备形状的第一段原文：

\begin{codebox}
\codeline{mfgsh:EquipmentShape}
\codeline{\hspace*{4\ttcw}a\ sh:NodeShape\ ;}
\codeline{\hspace*{4\ttcw}sh:targetClass\ mfg:Equipment\ ;}
\codeblank{}
\codeline{\hspace*{4\ttcw}\textcolor{cmtgray}{\#\ 序列号：必填、唯一格式\ EQ-YYYY-NNNN}}
\codeline{\hspace*{4\ttcw}sh:property\ [}
\codeline{\hspace*{8\ttcw}sh:path\ mfg:hasSerialNumber\ ;}
\codeline{\hspace*{8\ttcw}sh:minCount\ 1\ ;}
\codeline{\hspace*{8\ttcw}sh:maxCount\ 1\ ;}
\codeline{\hspace*{8\ttcw}sh:datatype\ xsd:string\ ;}
\codeline{\hspace*{8\ttcw}sh:pattern\ "\^{}EQ-[0-9]\{4\}-[0-9]\{4\}\$"\ ;}
\codeline{\hspace*{8\ttcw}sh:message\ "设备必须有且仅有一个格式为\ EQ-YYYY-NNNN\ 的序列号"\ ;}
\codeline{\hspace*{4\ttcw}]\ ;}
\end{codebox}
\color{black}

\snipsource{ch07-knowledge-graph/examples/kg-quality-shacl.ttl}

逐行解剖。\texttt{sh:targetClass} 是锚：图谱中每个 \texttt{mfg:Equipment}
实例都会被这个形状检查，且沿子类闭包生效——\texttt{mfg:CNCLathe} 的实例同样在列，
7.3 节「归类越具体越好」的建议不会让设备逃出质检范围。
\texttt{sh:path} 指明受检属性；
\texttt{sh:minCount 1} 与 \texttt{sh:maxCount 1} 合起来读作「有且仅有一个」——
前者是封闭语义的必填检查（正是 OWL 给不了的那种），后者掐灭重复值；
\texttt{sh:datatype} 锁定字面量类型，7.3 节「字面量必须带类型」的纪律在此交由机器执行；
\texttt{sh:pattern} 用正则锁死序列号格式 \texttt{EQ-YYYY-NNNN}，错一位、少一段都过不去；
末尾的 \texttt{sh:message} 是纯运营字段：校验引擎会把这句话原样写进违规报告，
它的读者是凌晨被告警叫起来的数据管理员——请写人话，写清「该是什么样」。

完整的设备形状还有三段属性形状（见仓库文件）。
功率必须是大于 0 且不超过 10000 千瓦的浮点数——
上界看起来多余，直到有人把单位填成瓦、一台「15500 千瓦」的车床来闯关那天：
7.5 节规格归一化漏掉的，这里来兜底。
状态必须取自受控词表——\texttt{sh:in} 一行列出
Running、Idle、Maintenance、Fault 四个词表个体，
7.3 节把代码表升格为受控词表的投资在此兑现第二次收益。
位置必须指向 \texttt{mfg:WorkCell} 的实例——\texttt{sh:class}
检查对象属性指向的类别，挂错类型的边过不了关。
维护记录形状则示范了\textbf{跨字段约束}：\texttt{sh:lessThanOrEquals}
声明维护日期不得晚于记录创建日期——两个字段之间的相对关系同样能进形状。
把常用约束家族排成速查表：

\begin{center}
\begin{tabularx}{\linewidth}{@{}p{0.18\linewidth}p{0.40\linewidth}X@{}}
\toprule
\textbf{约束家族} & \textbf{代表词汇} & \textbf{把关内容} \\
\midrule
基数 & \texttt{sh:minCount}、\texttt{sh:maxCount} & 必填与唯一 \\
取值类型 & \texttt{sh:datatype}、\texttt{sh:class} & 字面量类型、对象所属类 \\
数值范围 & \texttt{sh:minExclusive}、\texttt{sh:maxInclusive} & 区间，如大于 0 且不超过 10000 \\
字符串格式 & \texttt{sh:pattern} & 编号、代码等格式正则 \\
枚举 & \texttt{sh:in} & 锁定受控词表取值 \\
字段间比较 & \texttt{sh:lessThanOrEquals} 等 & 跨属性的相对关系 \\
自定义 & \texttt{sh:sparql} & 任意可查询的业务规则 \\
\bottomrule
\end{tabularx}
\end{center}

\subsection{SPARQL 约束：表达力的最后一公里}

核心约束词汇能覆盖大多数场景，剩下的复杂业务规则交给 \texttt{sh:sparql}。
看排产一致性的例子：

\begin{codebox}
\codeline{\textcolor{cmtgray}{\#\ ==============================}}
\codeline{\textcolor{cmtgray}{\#\ 3.\ 任务排产一致性形状（SPARQL约束）}}
\codeline{\textcolor{cmtgray}{\#\ ==============================}}
\codeline{\textcolor{cmtgray}{\#\ 复杂业务规则用\ sh:sparql\ 表达：任务结束时间必须晚于开始时间}}
\codeline{mfgsh:TaskIntervalShape}
\codeline{\hspace*{4\ttcw}a\ sh:NodeShape\ ;}
\codeline{\hspace*{4\ttcw}sh:targetClass\ mfg:Task\ ;}
\codeline{\hspace*{4\ttcw}sh:sparql\ [}
\codeline{\hspace*{8\ttcw}a\ sh:SPARQLConstraint\ ;}
\codeline{\hspace*{8\ttcw}sh:message\ "任务结束时间必须晚于开始时间"\ ;}
\codeline{\hspace*{8\ttcw}sh:prefixes\ mfg:\ ;}
\codeline{\hspace*{8\ttcw}sh:select\ """}
\codeline{\hspace*{12\ttcw}SELECT\ \$this\ WHERE\ \{}
\codeline{\hspace*{16\ttcw}\$this\ mfg:startTime\ ?s\ .}
\codeline{\hspace*{16\ttcw}\$this\ mfg:endTime\ \ \ ?e\ .}
\codeline{\hspace*{16\ttcw}FILTER\ (?e\ <=\ ?s)}
\codeline{\hspace*{12\ttcw}\}}
\codeline{\hspace*{8\ttcw}"""\ ;}
\codeline{\hspace*{4\ttcw}]\ .}
\codeblank{}
\end{codebox}
\color{black}

\snipsource{ch07-knowledge-graph/examples/kg-quality-shacl.ttl}

这种写法的思维方向是反的——不是「定义什么是好的」，而是「查出什么是坏的」。
\texttt{sh:select} 里是一段标准 SPARQL：
变量 \texttt{\$this} 由校验引擎依次绑定到目标类 \texttt{mfg:Task} 的每个实例，
查询体取出任务的开始与结束时间，
\texttt{FILTER} 专门筛出「结束时间竟不晚于开始时间」的反例——
查到谁，谁就是违规者，\texttt{sh:message} 随之写进报告。
理论上，SPARQL 约束把 SHACL 的表达力推到「凡 SPARQL 能查的都能校验」：
跨节点、聚合统计、多跳路径统统可以入约束。
工程上请克制：它可读性差（「写出坏例查询」的双重否定费脑），
不同引擎的支持程度参差，性能上等于每批数据多跑一次完整查询。
守则：\textbf{能用核心约束表达的，绝不动用 SPARQL 约束}；
确实动用的，把业务含义老老实实写进 \texttt{sh:message} 与注释。

\subsection{severity 分级：让闸门活下来}

形状写好了，还剩一个运营问题：是不是所有违规都该一律拦下？

\begin{codebox}
\codeline{\textcolor{cmtgray}{\#\ ==============================}}
\codeline{\textcolor{cmtgray}{\#\ 4.\ 入库闸门的使用约定}}
\codeline{\textcolor{cmtgray}{\#\ ==============================}}
\codeline{\textcolor{cmtgray}{\#\ 校验级别约定（在流水线中区分处理，见\ kg-construction-pipeline.txt）：}}
\codeline{\textcolor{cmtgray}{\#\ \ \ sh:Violation\ \ ——\ 拒绝入库，进隔离区}}
\codeline{\textcolor{cmtgray}{\#\ \ \ sh:Warning\ \ \ \ ——\ 允许入库，生成数据质量工单}}
\codeline{\textcolor{cmtgray}{\#\ 本文件形状默认\ sh:Violation；}}
\codeline{\textcolor{cmtgray}{\#\ 如需降级为警告，在对应\ PropertyShape\ 中加：}}
\codeline{\textcolor{cmtgray}{\#\ \ \ sh:severity\ sh:Warning\ ;}}
\end{codebox}

\snipsource{ch07-knowledge-graph/examples/kg-quality-shacl.ttl}

一刀切的闸门活不长。所有违规统统硬拦，上游一次小改版就会让整批数据滞留隔离区，
业务等不起，最后的结局往往是有人「临时」绕过闸门上线——闸门从此名存实亡。
SHACL 用 \texttt{sh:severity} 给违规分级，仓库文件约定了两档的管道语义：
\texttt{sh:Violation} 拒绝入库、进隔离区、触发告警；
\texttt{sh:Warning} 放行入库、生成数据质量工单、限期消化。
分级线怎么画？问一个问题：这条数据进库后，会不会污染主干查询与统计？
会的——序列号缺失（设备身份不明）、词表外状态（分组统计失真）、
范围外数值（平均值被拉飞）——必须 Violation；
不会的——备注缺失、非关键字段的格式瑕疵——降为 Warning。
两条运营纪律随之而来：
\textbf{Violation 要少而铁}，每一条都拦得理直气壮，业务方挑战时答得上来；
\textbf{Warning 要有人消化}，质量工单长期无人处理，团队就会学会无视它，
告警疲劳之后，Warning 等于没有。
隔离区的裁决沿用 7.1 节的三选一（修数据、修管道、改本体），
每次裁决记录在案——这些记录是 7.9 节漂移检测最值钱的信号源。

\subsection{形状进 CI：校验即测试}

形状文件的最佳类比是\textbf{单元测试套件}：
它和代码一样进版本库、走评审、随被测对象协同演化。
仓库文件头部给了最小运行方式——\texttt{pip install pyshacl} 之后一行命令：

\begin{center}
\texttt{pyshacl -s kg-quality-shacl.ttl -d instance-data.ttl}
\end{center}

校验在两个位置执行。\textbf{管道运行时}：每批候选三元组过闸门逐批校验，
这是 7.1 节定下的主战场。\textbf{持续集成（CI）时}：
本体变更、映射规则变更或形状自身变更的每一次提交，
都对一份金标准样例数据跑全量校验——
映射改坏了字面量类型、本体类改名导致形状失配，
都会在合并之前亮红灯，而不是在生产隔离区爆仓之后才被发现。
形状还反过来服务本体演化：本体打算收紧某条公理时，
先把对应的新形状对存量数据跑一遍，
违规数量就是未来迁移工作量最诚实的预算（7.9 节展开）。
第 9 章会把「先校验、后写入」的闸门写成可单步调试的代码骨架，
届时本节这份形状文件将原样上岗。

\section{存储与查询架构}

闸门之后，三元组得有个家。选库讨论会上常见两派：
标准派主张三元组库——「本体生态全在这边」；
开发派主张属性图数据库——「图数据库不就是存图的吗，社区大、上手快」。
两边都只说对了一半。本节先把两种数据模型的真实差异摆上桌面，
再用同一个业务问题对照两种问法，
最后给出选型判据与一条兼得的混合架构。

\subsection{两种数据模型}

\begin{codebox}
\codeline{\textcolor{cmtgray}{\#\ ==============================}}
\codeline{\textcolor{cmtgray}{\#\ 1.\ 两种数据模型}}
\codeline{\textcolor{cmtgray}{\#\ ==============================}}
\codeline{\textcolor{cmtgray}{\#\ RDF三元组模型：}}
\codeline{\textcolor{cmtgray}{\#\ \ \ (mfg:Lathe\_003,\ mfg:locatedIn,\ mfg:WorkCell\_A)}}
\codeline{\textcolor{cmtgray}{\#\ \ \ 一切皆三元组；边本身不带属性}}
\codeline{\textcolor{cmtgray}{\#\ \ \ （边属性需具名图或RDF-star：<<:s\ :p\ :o>>\ :since\ "2026-01"\ ）}}
\codeblank{}
\codeline{\textcolor{cmtgray}{\#\ 属性图模型（Labeled\ Property\ Graph）：}}
\codeline{\textcolor{cmtgray}{\#\ \ \ 节点\ (:Equipment\ \{name:'Lathe\_003',\ power:15.5\})}}
\codeline{\textcolor{cmtgray}{\#\ \ \ 边\ \ \ [:LOCATED\_IN\ \{since:'2026-01'\}]}}
\codeline{\textcolor{cmtgray}{\#\ \ \ 节点和边都可携带属性}}
\codeblank{}
\end{codebox}

\snipsource{ch07-knowledge-graph/examples/kg-storage-query.txt}

差异比「都叫图」深得多。RDF 把世界拆到原子粒度：一切皆三元组，
节点没有内部结构，连「属性」也是一条边；
而\textbf{边自身不携带属性}——想给「设备位于 A 区」补一句「自 2026 年 1 月起」，
三元组里放不下第四个位置。出路有三条：
把元数据挂到具名图上（7.9 节）；
用 RDF-star 扩展语法，把一条三元组整体当作主语再挂断言（文件注释给了示例）；
或者沿用第 4 章的老手艺，把关系升格为中介节点——
7.4 节的故障事件正是这么建的。
\textbf{属性图}（Labeled Property Graph，LPG）则把属性内置：
节点带键值对，边也带键值对，\texttt{since} 一个键就解决问题。

表面上 LPG 赢了便利，但 RDF 的「不便」换来的东西要看清楚。
其一，IRI 全球标识：跨系统合并图谱不会撞名，7.6 节的联邦场景全靠它。
其二，词汇与数据同一模型：本体本身就是一段 RDF，与实例同库共存——
7.1 节说「本体随图谱入库并继续生效」，这件事在 LPG 里没有对应物：
标签 \texttt{:Equipment} 只是一个字符串，没有公理、没有层次、没有定义。
其三，标准栈互操作：SPARQL、OWL、SHACL 围绕同一个数据模型生长，
前两节的全部机制对三元组库开箱即用。
LPG 的属性语义则由应用约定：\texttt{status:'Idle'} 里的 \texttt{'Idle'}
能取哪些值、什么含义，数据库不知道，知道的是你团队的文档——
而文档拦不住一条拼错的状态。

\subsection{同一个问题的两种问法}

需求：找出 A 区所有空闲设备的名称与功率。两种写法并排：

\begin{codebox}
\codeline{\textcolor{cmtgray}{\#\ SPARQL（三元组库）：}}
\codeline{\textcolor{cmtgray}{\#\ \ \ PREFIX\ mfg:\ <http://example.org/manufacturing\#>}}
\codeline{\textcolor{cmtgray}{\#\ \ \ SELECT\ ?name\ ?power\ WHERE\ \{}}
\codeline{\textcolor{cmtgray}{\#\ \ \ \ \ \ \ ?e\ \ rdf:type/rdfs:subClassOf*\ mfg:Equipment\ ;}}
\codeline{\textcolor{cmtgray}{\#\ \ \ \ \ \ \ \ \ \ \ mfg:hasStatus\ \ mfg:Idle\ ;}}
\codeline{\textcolor{cmtgray}{\#\ \ \ \ \ \ \ \ \ \ \ mfg:locatedIn\ \ mfg:WorkCell\_A\ ;}}
\codeline{\textcolor{cmtgray}{\#\ \ \ \ \ \ \ \ \ \ \ mfg:hasName\ \ \ \ ?name\ ;}}
\codeline{\textcolor{cmtgray}{\#\ \ \ \ \ \ \ \ \ \ \ mfg:hasPower\ \ \ ?power\ .}}
\codeline{\textcolor{cmtgray}{\#\ \ \ \}}}
\codeblank{}
\codeline{\textcolor{cmtgray}{\#\ Cypher（Neo4j）：}}
\codeline{\textcolor{cmtgray}{\#\ \ \ MATCH\ (e:Equipment\ \{status:'Idle'\})-[:LOCATED\_IN]->(w:WorkCell\ \{name:'A'\})}}
\codeline{\textcolor{cmtgray}{\#\ \ \ RETURN\ e.name,\ e.power}}
\end{codebox}

\snipsource{ch07-knowledge-graph/examples/kg-storage-query.txt}

先读 SPARQL 版。最关键的是第一行属性路径
\texttt{rdf:type/rdfs:subClassOf*}：它沿子类闭包匹配，
个体不论被归入 \texttt{mfg:Equipment} 还是它的任何子类（数控车床、加工中心），
一网打尽——这一行就是 7.1 节「查询时沿子类闭包展开」的落点，
也是「本体在查询里生效」的最短证明；
若写死成 \texttt{rdf:type mfg:Equipment}，
7.3 节鼓励的「归类越具体越好」反而会让设备从结果里消失。
其后几行用分号续写同一主语：状态等于词表个体 \texttt{mfg:Idle}（不是字符串），
位置等于 \texttt{mfg:WorkCell\_A}，再取名称与功率两个变量返回。

再读 Cypher 版。\texttt{MATCH} 后面是一段「ASCII 画图」：
节点写在圆括号里、关系写在方括号里、箭头标出方向，模式即查询，
两行解决问题，紧凑得近乎漂亮——这是开发者喜欢它的全部理由。
紧凑的代价藏在两处。
其一，\textbf{没有子类闭包}：\texttt{:Equipment} 标签匹配不到
只打了 \texttt{:CNCLathe} 标签的节点；
要么入库时给每个节点把祖先标签全部打上（把推理前移成建模纪律），
要么查询时手工枚举所有子类标签。
其二，\texttt{'Idle'} 是\textbf{裸字符串}：拼成小写不会报错，
只会静默返回空集；三元组库里词表个体写错则 IRI 根本不存在，错误暴露得更早。

文件里还有一组多跳对照：「产品 P001 延期会影响哪些订单」。
SPARQL 要逐跳写明三元组模式；Cypher 把多跳串进一个图样式。
变长路径的差异更有意思：SPARQL 属性路径 \texttt{mfg:precedes+}
取传递闭包，但拿不回中间路径经过了谁；
Cypher 的 \texttt{[:PRECEDES*1..5]} 既能限定跳数区间、又能把路径本身返回——
做追溯链可视化、瓶颈分析时顺手得多。这个差异会在选型表里再次现身。

\subsection{选型决策表}

\begin{codebox}
\codeline{\textcolor{cmtgray}{\#\ ==============================}}
\codeline{\textcolor{cmtgray}{\#\ 3.\ 投影\ adapter\ 选型决策表}}
\codeline{\textcolor{cmtgray}{\#\ ==============================}}
\codeline{\textcolor{cmtgray}{\#\ 维度\ \ \ \ \ \ \ \ \ \ \ \ \ \ 三元组库投影\ \ \ \ \ \ \ \ \ \ \ \ \ \ 属性图投影}}
\codeline{\textcolor{cmtgray}{\#\ ----------------\ \ --------------------------\ \ --------------------------}}
\codeline{\textcolor{cmtgray}{\#\ 语义交换保真\ \ \ \ \ \ RDF/OWL/SHACL保真度高\ \ \ \ \ \ \ 需显式映射公理与闭包}}
\codeline{\textcolor{cmtgray}{\#\ 标准化\ \ \ \ \ \ \ \ \ \ \ \ SPARQL/W3C，adapter易替换\ \ \ \ \ Cypher/GQL方言仍有差异}}
\codeline{\textcolor{cmtgray}{\#\ 边属性\ \ \ \ \ \ \ \ \ \ \ \ 具名图/RDF-star表达\ \ \ \ \ \ \ \ \ 原生支持}}
\codeline{\textcolor{cmtgray}{\#\ 图算法\ \ \ \ \ \ \ \ \ \ \ \ 较弱\ \ \ \ \ \ \ \ \ \ \ \ \ \ \ \ \ \ \ \ \ \ \ \ 强（路径/社区/中心性）}}
\codeline{\textcolor{cmtgray}{\#\ 联邦查询\ \ \ \ \ \ \ \ \ \ SERVICE关键字原生支持\ \ \ \ \ \ \ \ 通常需插件}}
\codeline{\textcolor{cmtgray}{\#\ 团队上手\ \ \ \ \ \ \ \ \ \ 需RDF知识\ \ \ \ \ \ \ \ \ \ \ \ \ \ \ \ \ \ \ 开发者友好}}
\codeblank{}
\codeline{\textcolor{cmtgray}{\#\ 投影选型经验：}}
\codeline{\textcolor{cmtgray}{\#\ \ \ 标准\ RDF\ 查询与交换优先\ \(\rightarrow\)\ 从\ Semantica\ receipt\ 生成只读三元组库投影}}
\codeline{\textcolor{cmtgray}{\#\ \ \ 路径分析/图算法密集（如供应链追溯）\(\rightarrow\)\ 生成只读属性图投影}}
\codeline{\textcolor{cmtgray}{\#\ \ \ 两头都重\ \(\rightarrow\)\ 从同一绑定\ snapshot\ 生成两类投影；两者都不得反向成为权威层}}
\codeblank{}
\end{codebox}
\color{black}

\snipsource{ch07-knowledge-graph/examples/kg-storage-query.txt}

六个维度逐行展开。
\textbf{本体与推理支持}是三元组库的主场：本体直接加载，
RDFS/OWL 推理与 SHACL 校验开箱即用；
属性图里没有公理的位置，子类闭包、属性传递都要在应用层重新发明一遍。
\textbf{标准化}决定可替换性：SPARQL 是 W3C 标准，
查询从 Jena 搬到 GraphDB 照跑，供应商绑定风险低；
属性图各家方言并立，GQL 标准化正在路上。
\textbf{边属性}与\textbf{图算法}是属性图的主场：
原生边属性免去建模绕路；路径搜索、社区发现、中心性等算法库成熟，
三元组库在这两项上要么繁琐、要么薄弱。
\textbf{联邦查询}：SPARQL 的 \texttt{SERVICE} 关键字原生支持跨端点查询，
7.6 节那种各自维护、声明关联的联邦格局天然受益。
\textbf{团队上手}：RDF 的概念门槛实打实存在，属性图对开发者更友好——
这一行不要轻视：工具再正确，团队用不起来就是不正确。

文件末尾的三行选型经验可以直接抄用：
推理与标准合规优先（本书主线场景）选三元组库；
路径分析与图算法密集（如供应链追溯）选属性图；两头都重，上混合架构。
落到操作上：把你最常问的十个业务问题写下来，
数一数几个依赖本体推理、几个依赖路径算法——
\textbf{选型选的不是「更好的库」，而是与工作负载匹配的库}。

\subsection{混合架构：权威层加分析投影}

两头都要的团队越来越多，混合架构的数据流值得画清楚：

\begin{center}
校验闸门 \(\rightarrow\) 三元组库（权威知识层） \(\rightarrow\) 投影管道
\(\rightarrow\) 属性图（分析视图） \(\rightarrow\) 图算法与可视化
\end{center}

三元组库是\textbf{单一事实源}：本体与全量实例都住在这里，
7.7 节的闸门只设在它的入口，溯源与具名图也只在这里维护。
属性图是\textbf{衍生品}：按分析需求投影出子图
（比如只取设备、工序与先后关系边），跑完算法可以整库丢弃重建，
\textbf{绝不回写}——一旦有人往分析库里直接写知识，两库开始漂移，
单一事实源失守，混合架构就退化成「两个互相怀疑的库」。
投影节奏二选一：周期性全量重建，实现简单、延迟以小时计；
订阅变更增量同步，接近实时、复杂度高出一截。
最后泼一盆冷水：双库不是免费午餐——两套运维、一条同步管道，
同步管道自己也要监控。默认从单库起步，
让真实出现的分析需求、而不是技术好奇心，来触发第二个库。

\subsection{把三元组库跑起来}

选型落锤，落地其实只要几条命令。仓库文件给了 Jena Fuseki 的最小部署：

\begin{codebox}
\codeline{\textcolor{cmtgray}{\#\ ==============================}}
\codeline{\textcolor{cmtgray}{\#\ 4.\ 三元组库部署示例（Jena\ Fuseki）}}
\codeline{\textcolor{cmtgray}{\#\ ==============================}}
\codeline{\textcolor{cmtgray}{\#\ 下载\ apache-jena-fuseki\ 后：}}
\codeline{\textcolor{cmtgray}{\#\ \ \ ./fuseki-server\ --loc=./tdb2-data\ /mfg\ \ \ \ \ \#\ TDB2持久化，端点\ /mfg}}
\codeline{\textcolor{cmtgray}{\#\ 加载本体与数据：}}
\codeline{\textcolor{cmtgray}{\#\ \ \ curl\ -X\ POST\ --data-binary\ @manufacturing.owl\ \textbackslash{}}}
\codeline{\textcolor{cmtgray}{\#\ \ \ \ \ \ \ \ -H\ 'Content-Type:\ application/rdf+xml'\ \textbackslash{}}}
\codeline{\textcolor{cmtgray}{\#\ \ \ \ \ \ \ \ http://localhost:3030/mfg/data}}
\codeline{\textcolor{cmtgray}{\#\ SPARQL端点：}}
\codeline{\textcolor{cmtgray}{\#\ \ \ 查询\ http://localhost:3030/mfg/query}}
\codeline{\textcolor{cmtgray}{\#\ \ \ 更新\ http://localhost:3030/mfg/update\ \ \ （生产环境应关闭或鉴权）}}
\codeblank{}
\codeline{\textcolor{cmtgray}{\#\ 开启OWL推理（assembler配置片段）：}}
\codeline{\textcolor{cmtgray}{\#\ \ \ :dataset\ a\ ja:RDFDataset\ ;}}
\codeline{\textcolor{cmtgray}{\#\ \ \ \ \ \ \ ja:defaultGraph\ [\ ja:reasoner}}
\codeline{\textcolor{cmtgray}{\#\ \ \ \ \ \ \ \ \ \ \ [\ ja:reasonerURL\ <http://jena.hpl.hp.com/2003/OWLFBRuleReasoner>\ ]\ ]\ .}}
\codeblank{}
\end{codebox}

\snipsource{ch07-knowledge-graph/examples/kg-storage-query.txt}

一行启动：\texttt{--loc} 指定 TDB2 持久化目录，\texttt{/mfg} 是数据集端点名。
第二步用 \texttt{curl} 把本体文件灌进去——注意灌的是 \texttt{manufacturing.owl}：
本体与实例同库共存，7.1 节的工作定义在部署层面坐实。
随后得到两个端点：查询端点 \url{http://localhost:3030/mfg/query} 开放给应用；
\textbf{更新端点在生产环境必须关闭或加鉴权}。
这不是一句普通的安全提醒：7.7 节的闸门设在管道里，
若更新端点对外裸奔，等于在闸门旁边开了一扇不设防的侧门，
任何一段脚本都能绕过全部校验直写库内，闸门形同虚设。
\textbf{全部写入走管道，库的写接口只对管道开放}——
这条纪律与 7.1 节「实例数据永远走管道」首尾呼应。
最后一段 assembler 配置展示如何给数据集挂接 OWL 推理器，
查询时自动展开推理结论。
第 5 章讨论过的推理代价在生产规模下更加尖锐：
全量 OWL 前向链在亿级三元组上开销高昂，
工程上常退守 RDFS 子集，或离线物化推理结论再入库；
物化结论的回收难题，7.6 节的 sameAs 事故已经领教过——
此处的对策同样是具名图：物化产物单独成图，重算时整图替换。
具名图在本章已经第三次救场，下一节给它一个正式的运维席位。

\section{图谱运维：让图谱长期可信}

很多团队把上线当成终点，而图谱的麻烦在于它必须\textbf{长期活着}：
数据每天进、本体定期改、业务悄悄变。
没有运维设计的图谱，半年后的典型症状是——
没人说得清某个值从哪来，没人敢删任何东西，新需求宁可另起炉灶。
本节把长期健康拆成四件事：
组织（具名图）、观测（监控）、演化（本体升级与数据迁移）、预警（漂移检测）。

\subsection{具名图：运维的最小操作单位}

第 5 章把具名图当作时态与溯源机制介绍；换到运维视角，
它升格为\textbf{图谱的目录结构}——没有它，亿级三元组是一锅粥；
有了它，每一类运维操作都有了抓手。仓库文件给出两种组织策略：

\begin{codebox}
\codeline{\textcolor{cmtgray}{\#\ ==============================}}
\codeline{\textcolor{cmtgray}{\#\ 5.\ 具名图组织策略}}
\codeline{\textcolor{cmtgray}{\#\ ==============================}}
\codeline{\textcolor{cmtgray}{\#\ 按来源分图：graph\ <http://example.org/graph/ledger>\ \ \ 台账}}
\codeline{\textcolor{cmtgray}{\#\ \ \ \ \ \ \ \ \ \ \ \ \ graph\ <http://example.org/graph/mes>\ \ \ \ \ \ MES}}
\codeline{\textcolor{cmtgray}{\#\ \ \ \ \ \ \ \ \ \ \ \ \ graph\ <http://example.org/graph/extracted>\ 文本抽取（低可信）}}
\codeline{\textcolor{cmtgray}{\#\ 按时间分图：graph\ <http://example.org/graph/20260321>\ \ 快照（见第5章时态）}}
\codeline{\textcolor{cmtgray}{\#\ 价值：来源可信度查询、按图回滚、抽取数据隔离审计}}
\end{codebox}

\snipsource{ch07-knowledge-graph/examples/kg-storage-query.txt}

\textbf{按来源分图}让前文的多处预告一次性落地。
7.4 节说「文本抽取的三元组进低可信具名图」，指的就是 \texttt{extracted} 图，
与台账图、MES 图物理隔离；
7.6 节冲突消解的第一优先「来源可信度」，前提正是「能按来源查」——
SPARQL 用 \texttt{GRAPH} 子句限定查询范围，
权威值到台账图里取，参考值才看抽取图；
审计抽取质量时，审计员只扫 \texttt{extracted} 一张图即可，不必大海捞针。
\textbf{按时间分图}承接第 5 章的时态机制：定期快照成图，
ABox 的历史可以按时间切片回溯，「上个季度这台设备归谁管」有据可查。

具名图还附赠一件运维利器——\textbf{按图回滚}。
某晚发现当天一批抽取数据系统性出错（比如提示词改坏了），
若它们混在默认图里，撤销要逐条甄别；
独立成图，则一条 \texttt{DROP GRAPH} 干净收回——
这是图谱世界里最接近「事务回滚」的武器。
图自身的 IRI 同样要规划：统一前缀 \texttt{http://example.org/graph/}、
名称稳定可预测，7.3 节的 IRI 铸造原则对图 IRI 一字不差地适用。

\subsection{监控：图谱的体检表}

监控的对象不只是库，而是\textbf{管道、库与人工队列的整体}——
图谱停止生长和长出垃圾同样是事故。一张最小可用的体检表：

\begin{center}
\begin{tabularx}{\linewidth}{@{}p{0.16\linewidth}XX@{}}
\toprule
\textbf{监控项} & \textbf{看什么} & \textbf{异常信号与第一反应} \\
\midrule
规模与增速 & 各具名图的三元组数与日增量 & 增速骤降疑似管道断流；骤增疑似上游异常或抽取失控 \\
闸门拒绝率 & 隔离区新增占批次的比例 & 持续上升首先怀疑上游字段改版或映射失配（7.3） \\
复核队列 & 对齐待裁决数与平均滞留时长 & 积压意味着新数据延迟入库，加人或调门限（7.5） \\
抽检精度 & 分关系类型的精确率与召回率趋势 & 无声下滑是文体漂移的先兆（7.4） \\
查询健康 & 关键业务查询的延迟 & 持续变慢提示数据膨胀，考虑物化或索引调整（7.8） \\
词汇覆盖 & 本体类与属性在实例中的使用分布 & 长期零实例的类，要么本体冗余、要么管道缺口（7.1） \\
\bottomrule
\end{tabularx}
\end{center}

每一行都直连前文某个环节，括号里给了回看的路标。
其中「词汇覆盖」最有本体味：它把镜头反过来对准 TBox——
本体里定义了的类在图谱中长期没有一个实例，
要么当年过度设计（回到第 3 章的范围控制反思一次），
要么某条该接的数据源根本没接通。
监控不只是保数据，也在替本体做体检。

\subsection{本体升级与 ABox 迁移}

仓库文件的版本策略一句话分好了工：
TBox 变更走发布流程（语义版本号加变更日志），
ABox 变更记录到具名图、支持时间切片回溯。
TBox 一侧的演化流程第 3 章讲过，这里补上图谱视角的关键难题：
\textbf{本体改了，库里百万级存量实例怎么办？}

按兼容性分两档。\textbf{向后兼容变更}——加类、加属性、放宽约束——
存量数据无需改动，发布即可，新数据按新本体进。
\textbf{破坏性变更}——类改名、一类拆成两类、公理收紧、词表删值——必须迁移，五步走：
其一，影响评估，用 SPARQL 数清受影响实例的规模与分布；
其二，编写迁移补丁，批量改写的更新脚本连同新旧本体一起评审；
其三，副本演练，在生产副本上跑通并核对前后计数；
其四，全量复检，7.7 节的形状此刻是迁移的验收标准——
新形状跑一遍迁移后的数据，违规清零才算完
（其实在动手之前，「新形状跑存量数据」的违规数就是迁移工作量最诚实的预算，
7.7 节末尾埋的这一手在此兑现）；
其五，切换并留底，旧版本数据保留为快照具名图——回得去，才敢往前走。
最后一条纪律：\textbf{本体发布与数据迁移必须在同一个变更窗口完成}，
否则图谱处于「新本体配旧数据」的撕裂态，闸门与查询两头报错。

\subsection{知识漂移检测}

设施都齐了，最后还剩一个安静的敌人：\textbf{知识漂移}——
图谱与它描述的现实渐行渐远。三种形态：
\textbf{现实漂移}，设备改造升级了，台账却没人更新——
CDC 只能捕获数据库里发生的变化，捕获不到「该录而没录」的变化，
图谱忠实地同步着一份过时的台账；
\textbf{文体漂移}，新班组、新缩写、新设备型号让维修日志的语言分布悄悄移动，
7.4 节预警过的抽取精度无声衰减；
\textbf{语义漂移}，业务术语的外延变了——「加工中心」悄悄涵盖了新购的五轴设备，
本体里的类定义还停在立项那年。

漂移没有专用传感器，但前文各环节早已埋好信号源：
\textbf{隔离区的裁决记录}——同一类裁决反复出现，
说明不是偶发脏数据而是系统性变化，7.1 节预告的回收在此兑现；
\textbf{抽检精度趋势}——持续下滑指向文体漂移（7.4 节）；
\textbf{词表外取值的频率}——同一个「非法状态值」反复被闸门拒绝，
往往意味着业务真的长出了新状态，该改的是词表而不是数据（7.7 节）；
\textbf{复核队列的案例构成}——某类设备的对齐存疑突然增多，
多半是命名习惯变了（7.5 节）。
响应按性质分流：数据错，修源头；管道错，修映射与提示词；
知识错，走本体变更流程——上一小节的迁移机制随即接棒。

至此回环闭合：运维信号驱动本体演化，本体演化经迁移落回图谱。
7.1 节工作定义里的第三个成分「治理」，到这里才算全部讲完——
它不是某个组件，而是从闸门到裁决、从监控到变更的这一整条回路。

\section{本章小结、思考题与文件指引}

\subsection{本章小结}

本章从「本体评审通过」走到「图谱长期可信」，
把一条完整的构建流水线铺在了制造车间里。
回头看，全部内容围绕四组关键词展开，恰好对应 7.1 节工作定义里的三个成分及其归宿。

\textbf{骨架与血肉}。TBox 与 ABox 之间横着一道规模断层：
骨架数百个概念靠评审保质量，血肉百万级三元组只能靠管道保质量——
实例数据永远走管道，手只用来写映射规则和裁决疑难案例，
这是全章第一条也是最硬的一条纪律。
流水线六个环节（接入、映射与抽取、对齐、闸门、入库、增量）各司其职，
本体校验闸门的三道检查——词汇、公理、形状——是贯穿始终的质量生命线，
未通过者进隔离区，由人按「修数据、修管道、改本体」三选一裁决。

\textbf{进料}。盘点先于接入：评估矩阵按知识价值、数据质量、获取成本、
更新频率给数据源排队，「先结构化主干、再非结构化抽取」的开工顺序
兼顾信任、锚点与风险三重论证。
映射式构建的四句口诀——行变个体、列变数据属性、外键列变对象属性、
代码表升格受控词表——配上「稳定主键铸 IRI、命名空间先行、IRI 绝不复用」的铸造纪律
与 CDC 增量补丁，让台账两周内长成图谱主干。
抽取式构建用本体类充当 NER 的标签集与关系抽取的目标模式，
故障这类多参与者事实铸成事件中介节点；
规则、NER 模型与 LLM 三代技术按综合处理成本混合编队，
LLM 的幻觉由「生成靠神经网络、把关靠符号系统」的闭环兜住；
抽取质量分槽位计算精确率与召回率，按关系类型分别报数，上线后持续抽检。

\textbf{融合与质量}。实体对齐先规则后模型：
编号正则、别名表、规格归一化承包工作量大头，
长尾交给名称、属性、关系三路加权相似度与双门限三段制——
机器吞掉高分低分两端，人只裁中间窄带，复核标注每月回流调参；
上线前随机抽 200 个合并决策审计，精确率不低于 98\% 才放行，
因为错合是会传播的污染型错误，代价远高于漏合。
表达对齐结论时牢记 \texttt{owl:sameAs} 的全量双向传播：
同一企业内部默认物理合并，跨系统联邦走 \texttt{skos:closeMatch} 升级阶梯，
强语义只授予经过验证的结论；合并之后的属性冲突，
按来源可信度、时效性、人工裁决三级消解，三者的共同前提是溯源。
SHACL 则以封闭检查补上 OWL 开放世界管不了的数据质量：
推理向外扩、校验向内收，能用核心约束就不动用 SPARQL 约束，
severity 分级让闸门在业务压力下活得下来，
形状文件像单元测试一样进版本库、进 CI、随本体协同演化。

\textbf{安身与长寿}。存储选型选的不是「更好的库」，而是与工作负载匹配的库：
本体推理与标准合规倾向三元组库，边属性与图算法密集倾向属性图，
两头都重就上「权威知识层加分析投影」的混合架构，且分析库绝不回写。
部署层面的纪律与开篇呼应：全部写入走管道，更新端点只对管道开放。
运维把长期健康拆成四件事：具名图按来源与时间组织三元组、提供按图回滚，
监控体检表盯住管道、库与人工队列的整体，
本体破坏性变更配五步 ABox 迁移并与发布同窗口完成，
漂移检测从隔离区裁决记录、抽检精度趋势、词表外取值频率与复核队列构成里
读出现实、文体、语义三种漂移——运维信号驱动本体演化，治理回路就此闭合。
这条回路没有终点：图谱的可信从来不是某个时点的状态，
而是被一天天运营出来的性质，停止运营之日就是信任开始流失之时。

给动手派读者一条练习路线：
先把四个仓库文件整篇通读一遍——正文为叙事顺畅省略的注释与细节都在里面；
然后装好 pyshacl，对照 7.7 节把形状文件跑在一份自己手造的实例数据上，
故意漏掉序列号、填一个 0 功率、写一个词表外状态，
亲眼看看三条违规报告分别长什么样；
再按 7.8 节的命令在本机启动 Fuseki，把本体与实例灌进同一个数据集，
用属性路径查询验证子类闭包确实生效；
最后回到 7.5 节，把「3号车床」对「CNC-Lathe-003」的 0.88 分亲手算一遍，
体会权重设计背后的领域判断，再把双门限上下挪动一档，
观察自动合并与人工复核的工作量怎样随之涨落。
四步走完，本章的流水线你就亲手摸过一遍了。

\subsection{思考题}

以下八道题按流水线顺序覆盖本章主干。
建议动笔写出三元组、形状或方案草图，而不只是口头作答；
多数题目需要把两个小节的结论拼起来用，动笔之前先回到对应小节把片段重读一遍。

\begin{noteblock}
\begin{enumerate}
\item \textbf{盘点排序}：你的车间多了一类数据源——设备传感器的时序读数（温度、振动）。
  用 7.2 节的评估矩阵给它逐列打分，论证它在接入计划里该排到什么位置、适合什么接入方式。
  提示：先想清楚海量原始点位该不该整体进图谱——哪一层是「知识」、哪一层只是「数据」。
\item \textbf{映射设计}：台账新增一列「负责班组」，取值是班组编号。
  按四句口诀判断它该映射成数据属性还是对象属性，写出对应的三元组形态并说明依据。
  提示：用正文那句话自测——这个值会不会被继续追问？班组自己有没有属性与关系要挂？
\item \textbf{IRI 事故推演}：某团队用设备名称铸 IRI，两年后全厂设备统一改名。
  列出图谱里至少三类会断裂的东西，再为「名称 IRI 迁移到主键 IRI」拟一套步骤。
  提示：迁移步骤可以倒过来抄 7.9 节的五步——影响评估那一步先用 SPARQL 数清受牵连的边。
\item \textbf{事件建模}：把维修日志改成扁平建模——症状、原因、处置直接挂在设备上。
  构造一个「两次故障互相搅浑」的具体反例，说明事件中介节点为什么不可省略。
  提示：让同一台车床先后因不同原因停机两次，再试着回答「哪次异响是轴承磨损引起的」。
\item \textbf{对齐手算}：沿用 7.5 节的例子，但这回两台设备功率对不上、所在工作单元也不同，
  属性相似度降到 0.4、关系相似度降到 0.3，名称相似度仍是 0.6。
  按权重算出总分并给出分流结果，再解释为什么「名字像」远不足以自动合并。
  提示：算完对照权重设计的初衷——名称恰恰是各系统分裂最严重、最不可信的那路证据。
\item \textbf{sameAs 清污}：7.6 节那条错误 sameAs 的推理结论已经物化入库。
  写一份清理预案：依赖哪些前置设施、按什么顺序撤销、用什么标准验证清理干净。
  提示：没有溯源与具名图这道题无解——先列出两者各自能提供哪条线索，再排操作顺序。
\item \textbf{形状编写}：为「维护记录」写一个节点形状，覆盖四条需求：
  必须关联一台设备、必须有维护日期、维护日期不晚于记录创建日期、备注缺失只警告不拦截。
  提示：对照约束家族速查表，四条需求恰好用到 \texttt{sh:class}、\texttt{sh:minCount}、
  \texttt{sh:lessThanOrEquals} 与 \texttt{sh:severity} 各一次。
\item \textbf{选型与运维}：业务方的十个高频问题里，六个要沿工艺路线做多跳追溯并展示路径，
  四个依赖设备分类的子类闭包统计。给出存储选型与数据流设计，
  并为上线后的第一张监控体检表挑出你认为最优先的三个指标。
  提示：先判断这是不是 7.8 节说的「两头都重」；指标按「图谱停止生长与长出垃圾同样是事故」去挑。
\end{enumerate}
\end{noteblock}

\subsection{本章文件指引}

本章全部代码片段节选自下列四个仓库文件。
与前几章略有不同，这四份文件本身就是一条流水线的设计文档：
按表中顺序通读，相当于把全章再走一遍——
正文为叙事取舍而未及展开的清单、注释与备选方案，都完整保留在文件里。

\begin{center}
\begin{tabularx}{\linewidth}{@{}>{\ttfamily\footnotesize}p{0.40\linewidth}X@{}}
\toprule
\normalfont\textbf{文件（ch07-knowledge-graph/examples/）} & \textbf{内容与阅读建议} \\
\midrule
kg-construction-pipeline.txt & 数据源盘点清单、流水线总览、映射式与抽取式构建的完整示例、校验闸门约定、增量更新与版本策略；配合 7.1 至 7.4 与 7.9 阅读，开篇与结尾的工程经验值得抄进项目白板 \\
entity-resolution.txt & 对齐问题的两种形态、规则归一化三件套、加权相似度与双门限阈值、sameAs 两种表达与陷阱、冲突消解优先级、抽样审计建议；配合 7.5 与 7.6 阅读，建议把相似度示例亲手验算 \\
kg-quality-shacl.ttl & 可直接运行的形状文件：设备完整性、维护时效性、排产一致性（SPARQL 约束）三组形状与入库闸门使用约定；装好 pyshacl 即可对自己的数据开跑，配合 7.7 阅读 \\
kg-storage-query.txt & 两种数据模型对照、同一查询的 SPARQL 与 Cypher 写法、六维选型决策表、Fuseki 部署命令、具名图组织策略；配合 7.8 与 7.9 阅读，多跳查询一组建议两种语法都亲手跑一遍 \\
\bottomrule
\end{tabularx}
\end{center}

至此，血肉长在了骨架上：本体提供词汇与公理，管道送来实例与关系，
治理回路让两者长期健康咬合——图谱开始回答业务问题了。
下一章把镜头转向正在改写整个领域格局的变量：大语言模型。
它既是迄今最高效的图谱构建工人，又是最需要图谱约束的「幻觉源」；
提示词里的本体词汇表、Text2SPARQL、GraphRAG 与幻觉控制将逐一登场。
本章修好的校验闸门，正是那场神经与符号协同的入场券。
